\ifdefined\gkver\else\def\gkver{student}\fi
\documentclass[\gkver]{gaokaozhenti}
\begin{document}

\chapter{1999年上海}

\begin{enumerate}
\item
%% number: 1
%% paperName: 1999年上海高考第 1 题 4 分
%% typeId: 1
%% type: 单选题
%% chapter: 光学
%% point: 光电效应
%% method: 无
%% score: 4
%% degree: 850
%% duplicateId: 0
%% body:
某单色光照射某金属时不能产生光电效应，则下述措施中可能使该金属产生光电效应的是\xzanswer{C}
\fourchoices[answer=C]
{延长光照时间}
{增大光的强度}
{换用波长较短的光照射}
{换用频率较低的光照射}
%% answer: C
%% memo:
\memoanswer{
本题原卷及材料中未提供详解，待补充。
}

\item
%% number: 2
%% paperName: 1999年上海高考第 2 题 4 分
%% typeId: 1
%% type: 单选题
%% chapter: 原子和原子核
%% point: 天然放射性
%% method: 无
%% score: 4
%% degree: 850
%% duplicateId: 0
%% body:
天然放射现象的发现揭示了\xzanswer{C}
\fourchoices[answer=C]
{原子不可再分}
{原子的核式结构}
{原子核还可再分}
{原子核由质子和中子组成}
%% answer: C
%% memo:
\memoanswer{
本题原卷及材料中未提供详解，待补充。
}

\item
%% number: 3
%% paperName: 1999年上海高考第 3 题 4 分
%% typeId: 1
%% type: 单选题
%% chapter: 交流电
%% point: 交流电
%% method: 无
%% score: 4
%% degree: 850
%% duplicateId: 0
%% body:
某交流发电机产生的感应电动势与时间的关系如图所示。如果其他条件不变，仅使线圈的转速加倍，则交流电动势的最大值和周期分别变为\xzanswer{B}
\onepicture[w=6cm]{\includesvg{figs/03}}
\fourchoices[answer=B]
{$400\UV$，$0.02\Us$}
{$200\UV$，$0.02\Us$}
{$400\UV$，$0.08\Us$}
{$200\UV$，$0.08\Us$}
%% answer: B
%% memo:
\memoanswer{
本题原卷及材料中未提供详解，待补充。
}

\item
%% number: 4
%% paperName: 1999年上海高考第 4 题 4 分
%% typeId: 1
%% type: 单选题
%% chapter: 直线运动
%% point: 竖直上下抛运动
%% method: 建立模型
%% score: 4
%% degree: 650
%% duplicateId: 0
%% body:
某同学身高 $1.8\Um$，在运动会上他参加跳高比赛，起跳后身体横着越过了 $1.8\Um$ 高度的横杆。据此可估算出他起跳时竖直向上的速度大约为（取 $g = 10\Umsq$）\xzanswer{B}
\fourchoices[answer=B]
{$2\Ums$}
{$4\Ums$}
{$6\Ums$}
{$8\Ums$}
%% answer: B
%% memo:
\memoanswer{
本题原卷及材料中未提供详解，待补充。
}

\item
%% number: 5
%% paperName: 1999年上海高考第 5 题 4 分
%% typeId: 1
%% type: 单选题
%% chapter: 机械波
%% point: 波的多解性
%% method: 无
%% score: 4
%% degree: 450
%% duplicateId: 0
%% body:
一列简谐横波向右传播，波速为 $v$。沿波传播方向上有相距为 $L$ 的 P、Q 两质点，如图所示。某时刻 P、Q 两质点都处于平衡位置，且 P、Q 间仅有一个波峰。经过时间 $t$，Q 质点第一次运动到波谷。则 $t$ 的可能值有\xzanswer{D}
\onepicture[w=5cm]{\includesvg{figs/05}}
\fourchoices[answer=D]
{1 个}
{2 个}
{3 个}
{4 个}
%% answer: D
%% memo:
\memoanswer{
本题原卷及材料中未提供详解，待补充。
}

\item
%% number: 6
%% paperName: 1999年上海高考第 6 题 4 分
%% typeId: 1
%% type: 单选题
%% chapter: 电磁感应
%% point: 楞次定律推广
%% method: 无
%% score: 4
%% degree: 450
%% duplicateId: 0
%% body:
如图（a）所示，竖直放置的螺线管与导线 abcd 构成回路，导线所围区域内有一垂直纸面向里的变化的匀强磁场，螺线管下方水平桌面上有一导体圆环。导线 abcd 所围区域内磁场的磁感强度按图（b）中哪一图线所表示的方式随时间变化时，导体圆环将受到向上的磁场作用力？\xzanswer{A}
\onepicture[w=12cm]{\includesvg{figs/06}}
%% answer: A
%% memo:
\memoanswer{
本题原卷及材料中未提供详解，待补充。
}

\item
%% number: 7
%% paperName: 1999年上海高考第 7 题 5 分
%% typeId: 2
%% type: 多选题
%% chapter: 曲线运动
%% point: 天体运动
%% method: 无
%% score: 5
%% degree: 650
%% duplicateId: 0
%% body:
把太阳系各行星的运动近似看作匀速圆周运动，则离太阳越远的行星\xzanswer{BCD}
\fourchoices[answer=BCD]
{周期越小}
{线速度越小}
{角速度越小}
{加速度越小}
%% answer: BCD
%% memo:
\memoanswer{
本题原卷及材料中未提供详解，待补充。
}

\item
%% number: 8
%% paperName: 1999年上海高考第 8 题 5 分
%% typeId: 2
%% type: 多选题
%% chapter: 光学
%% point: 光的折射
%% method: 无
%% score: 5
%% degree: 650
%% duplicateId: 0
%% body:
一束复色可见光射到置于空气中的平板玻璃上，穿过玻璃后从下表面射出，变为 a、b 两束平行单色光，如图所示。对于两束单色光来说\xzanswer{AD}
\onepicture[w=4.5cm]{\includesvg{figs/08}}
\fourchoices[answer=AD]
{玻璃对 a 光的折射率较大}
{a 光在玻璃中传播的速度较大}
{b 光每个光子的能量较大}
{b 光的波长较长}
%% answer: AD
%% memo:
\memoanswer{
本题原卷及材料中未提供详解，待补充。
}

\item
%% number: 9
%% paperName: 1999年上海高考第 9 题 5 分
%% typeId: 2
%% type: 多选题
%% chapter: 气体性质
%% point: 气体图象
%% method: 无
%% score: 5
%% degree: 450
%% duplicateId: 0
%% body:
一定质量的理想气体自状态 A 经状态 C 变化到状态 B 这一过程在 $V-T$ 图上表示如右图。则\xzanswer{AC}
\onepicture[w=5cm]{\includesvg{figs/09}}
\fourchoices[answer=AC]
{在过程 AC 中，外界对气体做功}
{在过程 CB 中，外界对气体做功}
{在过程 AC 中，气体压强不断变大}
{在过程 CB 中，气体压强不断变小}
%% answer: AC
%% memo:
\memoanswer{
本题原卷及材料中未提供详解，待补充。
}

\item
%% number: 10
%% paperName: 1999年上海高考第 10 题 5 分
%% typeId: 2
%% type: 多选题
%% chapter: 电场
%% point: 电势能电势差电场力做功
%% method: 无
%% score: 5
%% degree: 450
%% duplicateId: 0
%% body:
图中 a、b 为竖直向上的电场线上的两点，一带电质点在 a 点由静止释放，沿电场线向上运动，到 b 点恰好速度为零。下列说法中正确的是\xzanswer{ABD}
\onepicture[h=4.5cm]{\includesvg{figs/10}}
\fourchoices[answer=ABD]
{带电质点在 a、b 两点所受的电场力都是竖直向上的}
{a 点的电势比 b 点的电势高}
{带电质点在 a 点的电势能比在 b 点的电势能小}
{a 点的电场强度比 b 点的电场强度大}
%% answer: ABD
%% memo:
\memoanswer{
本题原卷及材料中未提供详解，待补充。
}

\item
%% number: 11
%% paperName: 1999年上海高考第 11 题 5 分
%% typeId: 2
%% type: 多选题
%% chapter: 运动和力
%% point: 瞬时加速度
%% method: 无
%% score: 5
%% degree: 350
%% duplicateId: 0
%% body:
如图所示，竖直光滑杆上套有一个小球和两根弹簧，两弹簧的一端各与小球相连，另一端分别用销钉 M、N 固定于杆上，小球处于静止状态。设拔去销钉 M 瞬间，小球加速度的大小为 $12\Umsq$。若不拔去销钉 M 而拔去销钉 N 瞬间，小球的加速度可能是（取 $g = 10\Umsq$）\xzanswer{BC}
\onepicture[h=4.5cm]{\includesvg{figs/11}}
\fourchoices[answer=BC]
{$22\Umsq$，竖直向上}
{$22\Umsq$，竖直向下}
{$2\Umsq$，竖直向上}
{$2\Umsq$，竖直向下}
%% answer: BC
%% memo:
\memoanswer{
本题原卷及材料中未提供详解，待补充。
}

\item
%% number: 12
%% paperName: 1999年上海高考第 12 题 4 分
%% typeId: 3
%% type: 填空题
%% chapter: 磁场
%% point: 安培力
%% method: 无
%% score: 4
%% degree: 650
%% duplicateId: 0
%% body:
在倾角为 $30^{\circ}$ 的光滑斜面上垂直纸面放置一根长为 $L$，质量为 $m$ 的直导体棒，一匀强磁场垂直于斜面向下，如图所示。当导体棒内通有垂直纸面向里的电流 $I$ 时，导体棒恰好静止在斜面上，则磁感强度的大小为 $B =$ \tkanswer[2.2cm]{$\frac{mg}{2IL}$}。
\onepicture[h=4cm, align=right]{\includesvg{figs/12}}
%% answer:
\jdanswer{
$\frac{mg}{2IL}$
}
%% memo:
\memoanswer{
本题原卷及材料中未提供详解，待补充。
}

\item
%% number: 13
%% paperName: 1999年上海高考第 13 题 4 分
%% typeId: 3
%% type: 填空题
%% chapter: 力物体的平衡
%% point: 力矩平衡
%% method: 无
%% score: 4
%% degree: 450
%% duplicateId: 0
%% body:
如图所示，质量不计的杆 O$_{1}$B 和 O$_{2}$A，长度均为 $l$，O$_{1}$ 和 O$_{2}$ 为光滑固定转轴，A 处有一凸起物搁在 O$_{1}$B 的中点，B 处用绳系在 O$_{2}$A 的中点，此时两短杆便组合成一根长杆。今在 O$_{1}$B 杆上的 C 点（C 为 AB 的中点）悬挂一重为 $G$ 的物体，则 A 处受到的支承力大小为 \tkanswer[1.4cm]{$\frac{G}{2}$}，B 处绳的拉力大小为 \tkanswer[1.2cm]{$G$}。
\onepicture[w=5cm, align=right]{\includesvg{figs/13}}
%% answer:
\jdanswer{
$\frac{G}{2}$，$G$
}
%% memo:
\memoanswer{
本题原卷及材料中未提供详解，待补充。
}

\item
%% number: 14
%% paperName: 1999年上海高考第 14 题 4 分
%% typeId: 3
%% type: 填空题
%% chapter: 光学
%% point: 光的电磁说
%% method: 近似与估算
%% score: 4
%% degree: 650
%% duplicateId: 0
%% body:
如图所示，古希腊某地理学家通过长期观测，发现 6 月 21 日正午时刻，在北半球 A 城阳光与铅直方向成 $7.5^{\circ}$ 角下射，而在 A 城正南方，与 A 城地面距离为 $L$ 的 B 城，阳光恰好沿铅直方向下射，射到地球的太阳光可视为平行光。据此他估算出了地球的半径。试写出估算地球半径的表达式 $R =$ \tkanswer[2.4cm]{$\frac{24L}{\pi}$}。
\onepicture[h=4.5cm, align=right]{\includesvg{figs/14}}
%% answer:
\jdanswer{
$\frac{24L}{\pi}$
}
%% memo:
\memoanswer{
本题原卷及材料中未提供详解，待补充。
}

\item
%% number: 15
%% paperName: 1999年上海高考第 15 题 4 分
%% typeId: 3
%% type: 填空题
%% chapter: 电路
%% point: 串并联电路
%% method: 等效替代
%% score: 4
%% degree: 650
%% duplicateId: 0
%% body:
如图所示电路由 8 个不同的电阻组成，已知 $R_{1} = 12\UO$，其余电阻阻值未知，测得 A、B 间的总电阻为 $4\UO$。今将 $R_{1}$ 换成 $6\UO$ 的电阻，则 A、B 间的总电阻变为 \tkanswer[1.2cm]{3} $\UO$。（提示：用等效替代法）
\onepicture[h=4cm, align=right]{\includesvg{figs/15}}
%% answer:
\jdanswer{
3
}
%% memo:
\memoanswer{
本题原卷及材料中未提供详解，待补充。
}

\item
%% number: 16
%% paperName: 1999年上海高考第 16 题 4 分
%% typeId: 3
%% type: 填空题
%% chapter: 直线运动
%% point: 匀速直线运动
%% method: 无
%% score: 4
%% degree: 450
%% duplicateId: 0
%% body:
天文观测表明，几乎所有远处的恒星（或星系）都在以各自的速度背离我们而运动，离我们越远的星体，背离我们运动的速度（称为退行速度）越大，也就是说，宇宙在膨胀，不同星体的退行速度 $v$ 和它们离我们的距离 $r$ 成正比，即

$v = Hr$

式中 $H$ 为一常量，称为哈勃常数，已由天文观察测定。为解释上述现象，有人提出一种理论，认为宇宙是从一个大爆炸的火球开始形成的。假设大爆炸后各星体即以不同的速度向外匀速运动，并设想我们就位于其中心，则速度越大的星体现在离我们越远。这一结果与上述天文观测一致。

由上述理论和天文观测结果，可估算宇宙年龄 $T$，其计算式为 $T =$ \tkanswer[2cm]{$\frac{1}{H}$}，根据近期观测，哈勃常数 $H = 3\times10^{-2}$ 米/秒·光年，其中光年是光在一年中行进的距离，由此估算宇宙的年龄约为 \tkanswer[2.4cm]{$1\times10^{10}$} 年。
%% answer:
\jdanswer{
$\frac{1}{H}$，$1\times10^{10}$
}
%% memo:
\memoanswer{
本题原卷及材料中未提供详解，待补充。
}

\item
%% number: 17
%% paperName: 1999年上海高考第 17 题 6 分
%% typeId: 4
%% type: 实验题
%% chapter: 机械振动
%% point: 单摆测重力加速度
%% method: 无
%% score: 6
%% degree: 650
%% duplicateId: 0
%% body:
在做“用单摆测定重力加速度”的实验时，用摆长 $l$ 和周期 $T$ 计算重力加速度的公式是 $g =$ \tkanswer[2.6cm]{$\frac{4\pi^{2}l}{T^{2}}$}。如果已知摆球直径为 $2.00\Ucm$，让刻度尺的零点对准摆线的悬点，摆线竖直下垂，如图（a）所示，那么单摆摆长是 \tkanswer[4cm]{$0.8740\Um$ 或 $87.40\Ucm$}。如果测定了 40 次全振动的时间如图（b）中秒表所示，那么秒表读数是 \tkanswer[1.4cm]{75.2} 秒，单摆的摆动周期是 \tkanswer[1.4cm]{1.88} 秒。
\onepicture[w=8cm, align=right]{figs/17.png}
%% answer:
\jdanswer{
$\frac{4\pi^{2}l}{T^{2}}$，$0.8740\Um$ 或 $87.40\Ucm$（$0.874\Um$ 或 $87.4\Ucm$ 不扣分），75.2，1.88
}
%% memo:
\memoanswer{
本题原卷及材料中未提供详解，待补充。
}

\item
%% number: 18
%% paperName: 1999年上海高考第 18 题 8 分
%% typeId: 4
%% type: 实验题
%% chapter: 电磁感应
%% point: 验证电磁感应实验
%% method: 无
%% score: 8
%% degree: 650
%% duplicateId: 0
%% body:
如图为“研究电磁感应现象”的实验装置。
\begin{enumerate}
	\item
	将图中所缺的导线补接完整。
	\item
	如果在闭合电键时发现灵敏电流计的指针向右偏了一下，那么合上电键后（　　）

	（A）将原线圈迅速插入副线圈时，电流计指针向右偏转一下

	（B）将原线圈插入副线圈后，电流计指针一直偏在零点右侧

	（C）原线圈插入副线圈后，将滑动变阻器触头迅速向左拉时，电流计指针向右偏转一下

	（D）原线圈插入副线圈后，将滑动变阻器触头迅速向左拉时，电流计指针向左偏转一下
\end{enumerate}
\onepicture[w=6cm, align=right]{figs/18.png}
%% answer:
\jdanswer{
\begin{enumerate}
	\item
	如图所示

	\onepicture[w=6cm]{figs/18_an.png}

	\item
	A，D

\end{enumerate}
}
%% memo:
\memoanswer{
本题原卷及材料中未提供详解，待补充。
}

\item
%% number: 19
%% paperName: 1999年上海高考第 19 题 4 分
%% typeId: 4
%% type: 实验题
%% chapter: 气体性质
%% point: 验证玻意耳定律
%% method: 无
%% score: 4
%% degree: 650
%% duplicateId: 0
%% body:
某同学做“验证玻意耳定律”实验时，将注射器竖直放置，测得的数据如下表所示。发现第 5 组数据中的 $pV$ 乘积值有较大偏差。如果读数和计算无误，那么造成此偏差的原因可能是 \tkanswer[1.4cm]{温度升高} 或 \tkanswer[1.4cm]{漏入气体}。

\begin{center}
\begin{tabular}{|c|c|c|c|c|c|}
\hline
实验次序 & 1 & 2 & 3 & 4 & 5 \\
\hline
$p$（$10^{5}\UPa$） & 1.21 & 1.06 & 0.93 & 0.80 & 0.66 \\
\hline
$V$（$\UmL$） & 33.2 & 37.8 & 43.8 & 50.4 & 69.2 \\
\hline
$pV$（$10^{5}\UPa\cdot\UmL$） & 40.2 & 40.1 & 40.7 & 40.3 & 45.7 \\
\hline
\end{tabular}
\end{center}
%% answer:
\jdanswer{
温度升高，漏入气体
}
%% memo:
\memoanswer{
本题原卷及材料中未提供详解，待补充。
}

\item
%% number: 20
%% paperName: 1999年上海高考第 20 题 5 分
%% typeId: 1
%% type: 单选题
%% chapter: 直线运动
%% point: 打点计时器公式
%% method: 无
%% score: 5
%% degree: 450
%% duplicateId: 0
%% body:
为了测定某辆轿车在平直路上起动时的加速度（轿车起动时的运动可近似看作匀加速运动），某人拍摄了一张在同一底片上多次曝光的照片（如图）。如果拍摄时每隔 $2\Us$ 曝光一次，轿车车身总长为 $4.5\Um$，那么这辆轿车的加速度约为\xzanswer{B}
\onepicture[w=12cm]{\includesvg{figs/20}}
\fourchoices[answer=B]
{$1\Umsq$}
{$2\Umsq$}
{$3\Umsq$}
{$4\Umsq$}
%% answer: B
%% memo:
\memoanswer{
本题原卷及材料中未提供详解，待补充。
}

\item
%% number: 21
%% paperName: 1999年上海高考第 21 题 7 分
%% typeId: 4
%% type: 实验题
%% chapter: 电路
%% point: 测电动势与内阻
%% method: 无
%% score: 7
%% degree: 650
%% duplicateId: 0
%% body:
现有一阻值为 $10.0\UO$ 的定值电阻、一个电键、若干根导线和一个电压表，该电压表表面上有刻度但无刻度值，要求设计一个能测定某电源内阻的实验方案。（已知电压表内阻很大，电压表量程大于电源电动势，电源内阻约为几欧）要求：
\begin{enumerate}
	\item
	在右边方框中画出实验电路图。
	\item
	简要写出完成接线后的实验步骤。
	\item
	写出用测得的量计算电源内阻的表达式 $r =$ \tkanswer[3cm]{$\frac{N_{1}-N_{2}}{N_{2}}R$}。
\end{enumerate}
%% answer:
\jdanswer*{
\begin{enumerate}
	\item
	如图所示

	\onepicture[w=4.5cm]{\includesvg{figs/21_an}}

	\item
	\begin{steps}
		\item
		断开电键，记下电压表偏转格数 $N_{1}$。
		\item
		合上电键，记下电压表偏转格数 $N_{2}$。
	\end{steps}

	\item
	$r = \frac{N_{1}-N_{2}}{N_{2}}R$

\end{enumerate}
}
%% memo:
\memoanswer{
本题原卷及材料中未提供详解，待补充。
}

\item
%% number: 22
%% paperName: 1999年上海高考第 22 题 10 分
%% typeId: 6
%% type: 计算题
%% chapter: 光学
%% point: 光的折射
%% method: 无
%% score: 10
%% degree: 650
%% duplicateId: 0
%% body:
光线以入射角 $i$ 从空气向折射率 $n = \sqrt{2}$ 的透明媒质表面。
\begin{enumerate}
	\item
	当入射角 $i = 45^{\circ}$ 时，求反射光线与折射光线间的夹角 $\theta$。
	\item
	当入射角 $i$ 为何值时，反射光线与折射光线间的夹角 $\theta = 90^{\circ}$？
\end{enumerate}
\onepicture[w=5cm, align=right]{\includesvg{figs/22}}
%% answer:
\jdanswer{
\begin{enumerate}
	\item
	$\theta = 105^{\circ}$

	\item
	$i = \arctan\sqrt{2}$

\end{enumerate}
}
%% memo:
\memoanswer{
\onepicture[align=right, w=5cm]{\includesvg{figs/22_an}}

（1）设折射角为 $r$，由折射定律

$\frac{\sin i}{\sin r} = n$ ①

$\sin r = \frac{\sin i}{n} = 0.5$

得 $r = 30^{\circ}$ ②

而 $i' = i = 45^{\circ}$ ③

$\therefore\theta = 180^{\circ} - 45^{\circ} - 30^{\circ} = 105^{\circ}$ ④

（2）此时 $i' + r = 90^{\circ}$ ⑤

$\sin r = \cos i$

代入折射定律 $\tan i = \sqrt{2}$ ⑥

$i = \arctan\sqrt{2}$ ⑦

评分标准：全题 10 分。（1）6 分，写出①式 2 分，得出结果④式 4 分（仅得出②、③各 1 分）。（2）4 分，写出⑤式 1 分，得出结果⑦式 3 分（仅得出⑥得 2 分）。
}

\item
%% number: 23
%% paperName: 1999年上海高考第 23 题 12 分
%% typeId: 6
%% type: 计算题
%% chapter: 气体性质
%% point: 气体多过程
%% method: 无
%% score: 12
%% degree: 450
%% duplicateId: 0
%% body:
如图均匀薄壁 U 形管，左管上端封闭，右管开口且足够长。管的横截面积为 $S$，内装密度为 $\rho$ 的液体。右管内有一质量为 $m$ 的活塞搁在固定卡口上，卡口与左管上端等高，活塞与管壁间无摩擦且不漏气。温度为 $T_{0}$ 时，左、右管内液面高度相等，两管内空气柱长度均为 $L$，压强均为大气压强 $p_{0}$。现使两边温度同时逐渐升高，求：
\begin{enumerate}
	\item
	温度升高到多少时，右管活塞开始离开卡口上升？
	\item
	温度升高到多少时，左管内液面下降 $h$？
\end{enumerate}
\onepicture[h=5cm, align=right]{\includesvg{figs/23}}
%% answer:
\jdanswer{
\begin{enumerate}
	\item
	$T_{1} = \left(1 + \frac{mg}{p_{0}S}\right)T_{0}$

	\item
	$T_{2} = \frac{\left(p_{0} + \frac{mg}{S} + 2\rho gh\right)(L + h)}{p_{0}L}T_{0}$

\end{enumerate}
}
%% memo:
\memoanswer{
（1）右管内气体为等容过程：$\frac{p_{0}}{T_{0}} = \frac{p_{1}}{T_{1}}$ ①

$p_{1} = p_{0} + \frac{mg}{S}$ ②

由①、②式得 $T_{1} = \left(1 + \frac{mg}{p_{0}S}\right)T_{0}$ ③

（2）对左管内气体列出状态方程 $\frac{p_{0}LS}{T_{0}} = \frac{p_{2}V_{2}}{T_{2}}$ ④

$p_{2} = p_{0} + \frac{mg}{S} + 2\rho gh$ ⑤

$V_{2} = (L + h)S$ ⑥

$T_{2} = \frac{\left(p_{0} + \frac{mg}{S} + 2\rho gh\right)(L + h)}{p_{0}L}T_{0}$ ⑦

评分标准：全题 12 分。（1）5 分，写出①、②式各 2 分，得出结果③式 1 分。（2）7 分，写出④、⑤式各 2 分，写出⑥式 1 分，得出结果⑦式 2 分。
}

\item
%% number: 24
%% paperName: 1999年上海高考第 24 题 14 分
%% typeId: 6
%% type: 计算题
%% chapter: 电磁感应
%% point: 电磁感应受力分析
%% method: 无
%% score: 14
%% degree: 450
%% duplicateId: 0
%% body:
如图所示，长为 $L$、电阻为 $r = 0.30\UO$、质量为 $m = 0.10\Ukg$ 的金属棒 CD 垂直跨搁在位于水平面上的两条平行光滑金属导轨上，两导轨间距也为 $L$，金属棒与导轨间接触良好，导轨电阻不计，导轨左端接有阻值 $R = 0.50\UO$ 的电阻。量程为 $0\sim3.0\UA$ 的电流表串接在一条导轨上，量程为 $0\sim1.0\UV$ 的电压表接在电阻 $R$ 的两端。垂直导轨平面的匀强磁场向下穿过平面。现以向右恒定外力 $F$ 使金属棒向右移动。当金属棒以 $v = 2.0\Ums$ 的速度在导轨平面上匀速滑动时，观察到电路中的一个电表正好满偏，而另一个电表未满偏。问：
\begin{enumerate}
	\item
	此满偏的电表是什么表？说明理由。
	\item
	拉动金属棒的外力 $F$ 多大？
	\item
	此时撤去外力 $F$，金属棒将逐渐慢下来，最终停止在导轨上。求从撤去外力到金属棒停止运动的过程中通过电阻 $R$ 的电量。
\end{enumerate}
\onepicture[w=7cm, align=right]{\includesvg{figs/24}}
%% answer:
\jdanswer{
\begin{enumerate}
	\item
	电压表满偏。理由是：若电流表满偏，回路中的电流应是 $I = 3.0\UA$，则电压表的示数应是 $U = IR = 1.5\UV$ 大于电压表量程；这不符合题意；若是电压表满偏，这时回路的电流是 $I = U/R = 2.0\UA$，说明电流表未满偏。

	\item
	$F = 1.6\UN$

	\item
	$q = 0.25\UC$

\end{enumerate}
}
%% memo:
\memoanswer{
（1）电压表满偏。若电流表满偏，则 $I = 3\UA$，$U = IR = 1.5\UV$，大于电压表量程。

（2）由功能关系 $Fv = I^{2}(R + r)$ ①

而 $I = \frac{U}{R}$ ②

$F = \frac{U^{2}(R + r)}{R^{2}v}$ ③

代入数据得 $F = 1.6\UN$ ④

（3）由动量定理：$m\Delta v = IBL\Delta t$ ⑤

两边求和：$m\Delta v_{1} + m\Delta v_{2} + \cdots = BLI_{1}\Delta t_{1} + BLI_{2}\Delta t_{2} + \cdots$ ⑥

即：$mv = BLq$ ⑦

由电磁感应定律：$E = BLv$ ⑧

$E = I(R + r)$ ⑨

由⑦⑧⑨解得 $q = \frac{mv^{2}}{I(R + r)}$ ⑩

代入数据得 $q = 0.25\UC$ （11）

评分标准：全题 14 分。（1）3 分，判定电压表满偏 2 分，说明理由 1 分。（2）5 分，写出①式 2 分，写出②式 1 分，得出结果④式 2 分（仅得出③式得 1 分）。（3）6 分，写出⑥、⑦、⑧、⑨式各 1 分，得出结果（11）式 2 分（仅得出⑩式得 1 分）。
}

\item
%% number: 25
%% paperName: 1999年上海高考第 25 题 15 分
%% typeId: 6
%% type: 计算题
%% chapter: 运动和力
%% point: 动量守恒定律
%% method: 无
%% score: 15
%% degree: 250
%% duplicateId: 0
%% body:
如图所示，一辆质量 $m = 2\Ukg$ 的平板车左端放有质量 $M = 3\Ukg$ 的小滑块，滑块与平板车之间的摩擦系数 $\mu = 0.4$。开始时平板车和滑块共同以 $v_{0} = 2\Ums$ 的速度在光滑水平面上向右运动，并与竖直墙壁发生碰撞，设碰撞时间极短且碰撞后平板车速度大小保持不变，但方向与原来相反，平板车足够长，以至滑块不会滑到平板车右端。求：
\begin{enumerate}
	\item
	平板车第一次与墙壁碰撞后向左运动的最大距离。
	\item
	平板车第二次与墙壁碰撞前瞬间的速度 $v$。
	\item
	为使滑块始终不会滑到平板车右端，平板车至少多长？
\end{enumerate}
\onepicture[w=7cm, align=right]{\includesvg{figs/25}}
%% answer:
\jdanswer{
\begin{enumerate}
	\item
	$s = 0.33\Um$

	\item
	$v = 0.4\Ums$

	\item
	$l = 0.833\Um$

\end{enumerate}
}
%% memo:
\memoanswer{
\onepicture[align=right, w=9cm]{\includesvg{figs/25_an}}

（1）设第一次碰墙壁后，平板车向左移动 $s$，速度变为零。由于体系总动量向右，平板车速度为零时，滑块还在向右滑行。

由动能定理：

$-\mu Mgs = 0 - \frac{1}{2}mv_{0}^{2}$ ①

$s = \frac{mv_{0}^{2}}{2\mu Mg}$ ②

代入数据得：$s = 0.33\Um$ ③

（2）假如平板车在第二次碰墙前还未和滑块相对静止，那么其速度的大小肯定还是 $2\Ums$，滑块的速度则大于 $2\Ums$，方向均向右，这样就违反动量守恒。所以平板车在第二次碰墙前肯定已和滑块具有共同速度 $v$，此即平板车碰墙前瞬间的速度。

$Mv_{0} - mv_{0} = (M + m)v$ ④

$v = \frac{M - m}{M + m}v_{0}$ ⑤

代入数据得：$v = 0.4\Ums$ ⑥

（3）平板车与墙壁第一次碰撞后到滑块与平板车又达到共同速度 $v$ 前的过程，可用图（a）（b）（c）表示，图（a）为平板车与墙碰撞后瞬间滑块与平板车的位置，图（b）为平板车到达最左端时两者的位置，图（c）为平板车与滑块再次达到共同速度时两者的位置。在此过程中滑块动能减少等于摩擦力对滑块所做功 $\mu Mgs'$，平板车动能减少等于摩擦力对平板车所做 $\mu Mgs''$（平板车从 B 到 A 再回到 B 的过程中摩擦力做功为零），其中 $s'$、$s''$ 分别为滑块和平板车的位移。滑块和平板车动能总减少为 $\mu Mgl_{1}$。其中 $l_{1} = s' + s''$ 为滑块相对平板车的位移。此后，平板车与墙壁发生多次碰撞，每次情况与此类似，最后停在墙边。设滑块相对平板车总位移为 $l$，则有：

$\frac{1}{2}(M + m)v_{0}^{2} = \mu Mgl$ ⑦

$l = \frac{(M + m)v_{0}^{2}}{2\mu Mg}$ ⑧

代入数据得 $l = 0.833\Um$ ⑨

$l$ 即为平板车的最短长度。

评分标准：全题 15 分。（1）5 分，写出①式 3 分，得出结果③式 2 分（仅得出②式得 1 分）。（2）5 分，写出④式 3 分，得出结果⑥式 2 分（仅得出⑤式得 1 分）。（3）5 分，写出⑦式 3 分，得出结果⑨式 2 分（仅得出⑧式得 1 分）。
}

\end{enumerate}

\end{document}
